% ============================================================================
%  Chapter 8. What the Metrics Selected
%
%  \input file. Requires booktabs, amsmath, natbib in the parent preamble.
%  Bibliography: accord/thesis/refs.bib
%
%  This chapter supersedes the withdrawn NeurIPS submission on metric selection
%  under low prevalence. The arithmetic is unchanged; the claim attached to it
%  is narrowed in one direction (no descriptive claim about field practice) and
%  strengthened in another (the empirical situation is worse than the paper
%  described).
%
%  Traceability: Chapter 7, Table \ref{tab:systems} and Table \ref{tab:sweep};
%  accord/CORRECTIONS.md. All confusion matrices in this chapter are
%  reconstructed from reported precision and recall and are validated by
%  reproducing the independently reported MCC values.
% ============================================================================

\chapter{What the Metrics Selected}
\label{ch:metrics}

\section{Why This Chapter Is Separate}
\label{sec:met-overview}

Chapter~\ref{ch:eval} established that the detector has no ranking signal. This
chapter asks a different question, one that survives the specific system and is
the only part of the project that was previously submitted for external review
as a standalone contribution: at the prevalence of this evaluation set, what do
the standard summary statistics actually select for?

The question is separable from the system because it is arithmetic. Given an
evaluation set with $P$ positives and $Q$ negatives, balanced accuracy,
prevalence-adjusted bias-adjusted kappa, and Matthews correlation each induce an
ordering over classifiers, and that ordering encodes an exchange rate between
missed positives and avoided false positives. The exchange rate is a property of
$P$ and $Q$, not of the systems. It can be computed before any system is run,
and at $P = 12$ against $Q = 241$ it is large enough to be worth stating in a
paper.

The chapter proceeds in three parts. Section~\ref{sec:met-closedform} derives
the exchange rate. Section~\ref{sec:met-systems} applies it to the systems of
Chapter~\ref{ch:eval} and shows that all three statistics prefer a system
recovering $6$ of $12$ contradictions over one recovering $11$.
Section~\ref{sec:met-sharpened} then does the thing that the earlier submission
did not: it asks whether the preference was \emph{wrong}, and concludes that it
was not, which makes the original framing incorrect and the underlying
prescription stronger.

\section{The Prior Submission and What Reviewers Said}
\label{sec:met-prior}

The material in Sections~\ref{sec:met-closedform} and~\ref{sec:met-systems} was
submitted to NeurIPS 2026 and withdrawn. The reviews are summarised here because
the revision this chapter represents is a response to them, and because two of
the three objections were correct.

\paragraph{Significance, scored 1.} The reviewer wrote that the paper ``argues
against a failure mode that does not seem to exist in practice.'' This objection
is substantially correct as applied to the submitted framing. The paper asserted
that the field selects models by these statistics under these conditions, and
offered no evidence for that assertion. Reviewer JDYj asked for the evidence
directly. It does not exist, it was not collected, and this chapter therefore
withdraws the descriptive claim entirely. What remains is prescriptive, and
Section~\ref{sec:met-prescription} states it in that form.

\paragraph{Originality, scored 1 and 2.} Both reviewers characterised the
contribution as ``largely a class imbalance problem.'' This is correct as far as
it goes and is why Section~\ref{sec:bg-imbalance} of Chapter~\ref{ch:background}
now treats that literature at length. It is also not a complete answer. The
imbalance literature is overwhelmingly concerned with \emph{fitting} under skew:
resampling \citep{chawla2002smote}, cost-sensitive decision rules
\citep{elkan2001cost}, loss reshaping
\citep{lin2017focal,cui2019classbalanced,cao2019ldam,menon2021logitadjust}. The
one result in that literature which speaks most directly to evaluation,
\citet{vandengoorbergh2022harm}, establishes that imbalance corrections harm
calibration and do not improve discrimination, and that whatever they achieve is
achievable by moving a threshold. That is a statement about what corrections do
to a model. It is not a statement about which model a summary statistic picks,
and the question of which model a summary statistic picks is what
Section~\ref{sec:met-closedform} answers.

\paragraph{A third objection, from a different reviewer, that this thesis
accepts in full.} Reviewer 1yJj observed that there is no value in reporting on
a contaminated set, because that is training on the test set. Correct. The
contaminated configuration is withdrawn from this thesis in its entirety
(Chapter~\ref{ch:candidates}, Section~\ref{sec:splits}), and no number in this
chapter derives from it.

\section{The Closed-Form Result}
\label{sec:met-closedform}

\subsection{Setup}

Let an evaluation set contain $N$ items, of which $P$ are positive and
$Q = N - P$ are negative. A classifier produces a confusion matrix
$(\mathrm{TP}, \mathrm{FN}, \mathrm{FP}, \mathrm{TN})$ with
$\mathrm{TP} + \mathrm{FN} = P$ and $\mathrm{FP} + \mathrm{TN} = Q$. On the
evaluation set of Chapter~\ref{ch:eval}, restricted to the non-ambiguous pairs,
\[
N = 253, \qquad P = 12, \qquad Q = 241, \qquad \pi = P/N = 4.74\% .
\]

\subsection{Balanced accuracy}

Balanced accuracy \citep{brodersen2010balanced} is the unweighted mean of
sensitivity and specificity,
\[
\mathrm{BA} \;=\; \frac{1}{2}\!\left(\frac{\mathrm{TP}}{P} +
\frac{\mathrm{TN}}{Q}\right).
\]
For two classifiers $A$ and $B$,
\[
\mathrm{BA}(B) - \mathrm{BA}(A) \;=\;
\frac{1}{2}\!\left(\frac{\Delta\mathrm{TP}}{P} +
\frac{\Delta\mathrm{TN}}{Q}\right),
\]
with $\Delta\mathrm{TP} = \mathrm{TP}_B - \mathrm{TP}_A$ and
$\Delta\mathrm{TN} = \mathrm{TN}_B - \mathrm{TN}_A$. Suppose $B$ catches fewer
positives, $\Delta\mathrm{TP} = -\delta$ with $\delta > 0$. Then $B$ is
preferred exactly when

\begin{equation}
\boxed{\;\frac{\Delta\mathrm{TN}}{\delta} \;>\; \frac{Q}{P}\;}
\label{eq:exchange}
\end{equation}

that is, when $B$ avoids more than $Q/P$ of $A$'s false positives for every true
positive it forgoes. At $P = 12$ and $Q = 241$,
\[
\frac{Q}{P} = \frac{241}{12} = 20.08 .
\]

Equation~\eqref{eq:exchange} is the whole result. It contains no model, no
dataset beyond its two counts, and no assumption. It says that selecting by
balanced accuracy on this evaluation set is equivalent to declaring, in advance
and silently, that one missed contradiction is worth twenty avoided false
alarms.

\subsection{Prevalence-adjusted bias-adjusted kappa}

PABAK \citep{byrt1993pabak} was introduced to remove the sensitivity of Cohen's
$\kappa$ \citep{cohen1960kappa} to marginal imbalance, the sensitivity
documented as paradoxical by \citet{feinstein1990paradoxes} and analysed by
\citet{cicchetti1990resolving}. Its definition is
\[
\mathrm{PABAK} \;=\; 2 p_o - 1, \qquad
p_o = \frac{\mathrm{TP} + \mathrm{TN}}{N},
\]
where $p_o$ is observed agreement. For a classifier evaluated against a gold
label, $p_o$ is raw accuracy, so
\[
\mathrm{PABAK} = 2\,\mathrm{accuracy} - 1 ,
\]
a strictly increasing affine transform of accuracy. Its exchange rate follows
immediately: $B$ is preferred over $A$ exactly when
$\Delta\mathrm{TP} + \Delta\mathrm{TN} > 0$, that is, when
\[
\frac{\Delta\mathrm{TN}}{\delta} > 1 .
\]

This deserves emphasis because it is the opposite of what the measure's name
suggests. PABAK removes kappa's prevalence sensitivity by discarding the
marginal information that produced it, and what remains is a rescaling of
accuracy, whose implicit exchange rate is one avoided false positive per
forgone true positive \emph{in raw counts}. At $\pi = 4.74\%$ that is
substantially more hostile to the minority class than balanced accuracy, not
less. A measure adopted for its robustness to prevalence is, at this prevalence,
the least robust of the three. This observation is consistent with
\citet{delgado2019kappa}, who argue on related grounds that kappa should be
avoided as a classification performance measure.

\subsection{Matthews correlation}

MCC \citep{matthews1975mcc} is
\[
\mathrm{MCC} = \frac{\mathrm{TP}\cdot\mathrm{TN} -
\mathrm{FP}\cdot\mathrm{FN}}
{\sqrt{(\mathrm{TP}{+}\mathrm{FP})(\mathrm{TP}{+}\mathrm{FN})
(\mathrm{TN}{+}\mathrm{FP})(\mathrm{TN}{+}\mathrm{FN})}} .
\]
Unlike the other two it is not affine in the confusion cells, so it admits no
constant exchange rate and the comparison between two classifiers must be
evaluated numerically. It is nonetheless the statistic most often recommended
for imbalanced binary evaluation \citep{chicco2020mcc,powers2011evaluation}, and
it is the statistic Chapter~\ref{ch:eval} reports, so its behaviour on this
evaluation set is the practically relevant case.
Section~\ref{sec:met-systems} evaluates it directly.

\section{The Systems}
\label{sec:met-systems}

\subsection{Reconstructed confusion matrices, and their validation}

Chapter~\ref{ch:eval}, Table~\ref{tab:systems}, reports precision, recall, and
MCC for four systems. The full confusion matrix for each follows from precision,
recall, $P$, and $Q$ by arithmetic, and the reconstruction can be validated:
MCC is not used in the reconstruction, so reproducing the independently reported
MCC from the reconstructed cells is a genuine check.
Table~\ref{tab:met-confusion} reports both.

\begin{table}[t]
\centering
\small
\begin{tabular}{lrrrrrr}
\toprule
\textbf{System} & \textbf{TP} & \textbf{FN} & \textbf{FP} & \textbf{TN} &
\textbf{MCC} & \textbf{MCC} \\
 & & & & & \textbf{(derived)} & \textbf{(reported)} \\
\midrule
GPT-4o, zero-shot          & 6  & 6 & 65  & 176 & $+0.109$ & $+0.109$ \\
GPT-4o, few-shot           & 7  & 5 & 134 & 107 & $+0.012$ & $+0.012$ \\
ACCORD, $\tau = 0.5$       & 11 & 1 & 223 & 18  & $-0.007$ & $-0.007$ \\
Llama 3 8B (weak labeller) & 11 & 1 & 231 & 10  & $-0.044$ & $-0.044$ \\
\bottomrule
\end{tabular}
\caption{Confusion matrices reconstructed from the precision and recall reported
in Chapter~\ref{ch:eval}, Table~\ref{tab:systems}, on the $253$ non-ambiguous
pairs ($P = 12$, $Q = 241$). The final two columns are the check: MCC computed
from the reconstructed cells agrees with the independently reported MCC to three
decimal places for all four systems. The ACCORD row is independently
corroborated by Chapter~\ref{ch:eval}'s reported accuracy of $0.115$, which is
$(11 + 18)/253$. Nine further systems appeared in the withdrawn submission's
comparison table. [TODO: recover the remaining nine rows from the submission's
result artifacts and report them here on the identical $253$-pair basis, or
state that they cannot be reproduced and restrict the chapter to four.]}
\label{tab:met-confusion}
\end{table}

\subsection{What the three statistics select}

\begin{table}[t]
\centering
\small
\begin{tabular}{lrrrr}
\toprule
\textbf{System} & \textbf{Recall} & \textbf{Balanced} & \textbf{PABAK} &
\textbf{MCC} \\
 & \textbf{(of 12)} & \textbf{accuracy} & & \\
\midrule
GPT-4o, zero-shot          & 6  & \textbf{0.615} & $\mathbf{+0.439}$ & $\mathbf{+0.109}$ \\
GPT-4o, few-shot           & 7  & 0.514          & $-0.099$          & $+0.012$ \\
ACCORD, $\tau = 0.5$       & 11 & 0.496          & $-0.771$          & $-0.007$ \\
Llama 3 8B (weak labeller) & 11 & 0.479          & $-0.834$          & $-0.044$ \\
\bottomrule
\end{tabular}
\caption{All three statistics induce the same ordering, and all three place the
system that recovers $6$ of $12$ contradictions above the two that recover $11$.
Balanced accuracy and PABAK are derived from the cells in
Table~\ref{tab:met-confusion}; MCC is as reported in Chapter~\ref{ch:eval}.
Balanced accuracy separates the top and bottom systems by $0.14$; PABAK, being
a rescaling of accuracy, separates them by $1.27$ on a $[-1, 1]$ scale, which is
the numerical signature of a measure dominated by the majority class.}
\label{tab:met-selection}
\end{table}

Table~\ref{tab:met-selection} is the result the earlier submission reported.
Applying Equation~\eqref{eq:exchange} to the decisive pair makes the mechanism
explicit. Moving from ACCORD to GPT-4o zero-shot forgoes $\delta = 5$ true
positives and gains $\Delta\mathrm{TN} = 176 - 18 = 158$ true negatives. The
threshold is $5 \times 20.08 = 100.4$, so the preference flips once at least
$101$ of ACCORD's false positives are avoided. GPT-4o avoids $158$ of them, and
the flip is therefore not marginal: it clears the threshold by more than half
again.

\paragraph{The claim as originally made.} Balanced accuracy, PABAK, and MCC all
select a system that misses half the contradictions over one that misses one in
twelve, and the exchange rate that produces this is available in closed form
before any experiment is run. That claim is correct and this chapter does not
retract it.

\section{The Sharpened Claim}
\label{sec:met-sharpened}

\subsection{The question the earlier submission did not ask}

The submitted framing carried an unstated premise: that a system recovering $11$
of $12$ is the better system, and that the metrics, in preferring the one
recovering $6$, were preferring the worse. The reviewer who wrote that the paper
argues against a failure mode that does not seem to exist in practice was, on
this point, closer to right than the paper was.

Chapter~\ref{ch:eval}'s four measurements supply the missing information, and
they invert the reading. Consider what the $11$ of $12$ consists of. ACCORD
predicts \emph{contradiction} for $234$ of the $253$ non-ambiguous pairs. At
thresholds $0.05$, $0.10$, and $0.15$ its predicted-positive rate is exactly
$1.000$, because no score in the evaluation set falls below $0.15$
(Chapter~\ref{ch:eval}, Table~\ref{tab:sweep}). Its AUROC is $0.460$ with a
Mann-Whitney $p$ of $0.64$ and a bootstrap $95\%$ interval of $[0.287, 0.633]$;
its fifteen highest-scored pairs contain zero true contradictions. Its recall of
$11/12$ is the recall of a near-constant positive predictor, and it is not
evidence of detection.

\subsection{The corrected statement}

\textbf{The metrics were not choosing a worse system over a better one. They were
choosing among a field in which three of the four systems have no discriminative
ability at all, and the system they chose was the only one that had any.}

The evidence for that reading is in the MCC column of
Table~\ref{tab:met-selection} and is independent of the exchange-rate argument.
ACCORD sits at $-0.007$, indistinguishable from chance. The weak labeller it was
distilled from sits at $-0.044$, which is worse than chance. GPT-4o zero-shot at
$+0.109$, with $8.45\%$ precision against a base rate of $4.74\%$, is the only
system in the comparison whose predictions carry information about the labels.
Balanced accuracy, PABAK, and MCC all put it first. On this evaluation set they
were right.

\subsection{Why this makes the argument stronger, not weaker}

Three consequences follow, and they are the substance of the chapter.

\paragraph{The arithmetic is unchanged and is now unconfounded.}
Equation~\eqref{eq:exchange} holds regardless of which system happens to be
better. The exchange rate of $20.08$ true negatives per true positive is a
property of $P = 12$ and $Q = 241$, and a paper that selects by balanced
accuracy on such a set accepts that rate whether or not its particular
comparison lands on the right side of it. The earlier submission tied a general
arithmetic point to a specific empirical claim about which system was better,
and the empirical claim was wrong. Untying them removes the only part of the
argument that could be falsified by a better model.

\paragraph{The empirical situation is worse than the submission described.} The
submission's implicit worry was that a good high-recall system might be passed
over. What Chapter~\ref{ch:eval} shows is that the high-recall systems were not
good; they were constant. The failure mode a low-prevalence evaluation must
guard against is therefore not primarily ``the metric prefers the wrong system''
but ``every system in the comparison, including the one the metric prefers, may
be near the no-skill line, and a metric that ranks them will still return a
ranking.'' A statistic that produces an ordering over four systems with MCCs of
$+0.109$, $+0.012$, $-0.007$, and $-0.044$ has produced an ordering over a set
in which three members are noise. Nothing in the statistic says so. The
statistics that do say so are the ones Chapter~\ref{ch:eval} reports alongside
it: a threshold-free ranking measure with a confidence interval, a
significance test on the score distributions, an average precision against the
base rate, and a top-$k$ count.

\paragraph{Recall in isolation is the trap, not the metric.} The original
framing treated $11/12$ as the quantity to be protected. Both this chapter and
Chapter~\ref{ch:eval}, Section~\ref{sec:errors-reading}, now state the opposite:
recall alone is the most misleading statistic available in this setting, because
a constant-positive predictor achieves recall $1.0$ by construction. The
appropriate defence of a high-recall system is not that recall is high; it is a
demonstration that the predicted-positive rate is meaningfully below one and
that the score ordering separates the classes. ACCORD fails both.

\section{The Prescription}
\label{sec:met-prescription}

This chapter makes no claim about how the field selects models. It makes four
recommendations, each of which is checkable by a reviewer in one line and none
of which requires agreement with the diagnosis above.

\begin{enumerate}
\item \textbf{Report $Q/P$ next to any balanced accuracy.} It is the exchange
      rate that measure imposes, by Equation~\eqref{eq:exchange}. On this
      evaluation set it is $20.08$. A reader who is content with that rate can
      proceed; a reader who is not now knows to object, and did not before.

\item \textbf{Do not report PABAK as a prevalence-robust measure.} It is
      $2\,\mathrm{accuracy} - 1$. If accuracy would be an inappropriate summary
      for the evaluation set, PABAK is inappropriate for the same reason and to
      the same degree, and its name conceals this.

\item \textbf{Prespecify the primary metric, and state the cost ratio it
      encodes.} TRIPOD+AI \citep{collins2024tripodai} requires performance
      measures to be prespecified for clinical prediction models, and the
      cost-sensitive framing \citep{elkan2001cost} makes explicit that a
      decision threshold is a statement about relative misclassification cost.
      Where the costs are genuinely asymmetric and known only approximately,
      decision curve analysis \citep{vickers2006dca} reports performance across
      a range of threshold probabilities rather than committing to one.

\item \textbf{Report at least one measure that can say ``no system here has
      signal.''} A ranking statistic returns a ranking whether or not the ranked
      objects differ. The measures that refuse are a threshold-free
      discrimination estimate with an interval, a test on the score
      distributions, average precision against the base rate, and a top-$k$
      count. Chapter~\ref{ch:eval} reports all four, and all four agree that
      three of the four systems in Table~\ref{tab:met-selection} are noise.
\end{enumerate}

\paragraph{A fifth recommendation this thesis cannot yet support.} Where the
positive class is heterogeneous in severity, the natural evaluation is
severity-weighted rather than count-based. Chapter~\ref{ch:eval},
Table~\ref{tab:danger}, records that the twelve confirmed contradictions carry
danger ratings distributed $(2, 4, 2, 1, 3)$ across the scale $1$ to $5$, so
four of twelve are rated $4$ or $5$. A severity-weighted recall would count
those four more heavily. [TODO: report, for each system in
Table~\ref{tab:met-confusion}, which of the twelve true contradictions it
recovered, broken down by danger rating, and recompute recall with severity
weights. The per-item recovery is not recorded in the present result artifacts,
so the weighted comparison cannot be made here.]

\section{Proposed Figures}
\label{sec:met-figures}

\begin{figure}[t]
\centering
\framebox[\linewidth]{\parbox[c][5.5cm][c]{0.92\linewidth}{\centering
\textit{[Figure asset 8.1: the exchange rate. X axis, forgone true positives}
\textit{$\delta$ from 0 to 12. Y axis, additional true negatives $\Delta$TN from}
\textit{0 to 241. Three indifference lines through the origin: PABAK at slope 1,}
\textit{balanced accuracy at slope $Q/P = 20.08$, and a shaded band for MCC}
\textit{traced numerically. Plot the four systems of Table 8.1 as points}
\textit{relative to ACCORD as the origin; GPT-4o zero-shot sits at}
\textit{$(\delta, \Delta$TN$) = (5, 158)$, above both lines.]}}}
\caption{What each statistic trades. The reader should notice that the
indifference slope is a property of the evaluation set's class counts alone, and
that PABAK's slope of $1$ makes it the most majority-dominated of the three
despite its name.}
\label{fig:met-exchange}
\end{figure}

\begin{figure}[t]
\centering
\framebox[\linewidth]{\parbox[c][5cm][c]{0.92\linewidth}{\centering
\textit{[Figure asset 8.2: two panels sharing a y axis of four systems. Left,}
\textit{recall out of 12: GPT-4o zs 6, GPT-4o fs 7, ACCORD 11, Llama 11.}
\textit{Right, MCC with a vertical line at 0 labelled ``chance'': $+0.109$,}
\textit{$+0.012$, $-0.007$, $-0.044$. Annotate the ACCORD and Llama bars in the}
\textit{left panel with their predicted-positive counts, 234/253 and 242/253.}
\textit{Caption box: ``the high-recall systems predict positive almost}
\textit{everywhere''.]}}}
\caption{The sharpened claim in one view. The two systems with the highest
recall achieve it by predicting the positive class for over ninety percent of
the evaluation set, and both sit at or below chance on Matthews correlation. The
metrics that prefer the low-recall system are not preferring a worse detector;
they are the only summaries in the comparison that register the difference.}
\label{fig:met-sharpened}
\end{figure}

\section{Limitations of This Chapter}
\label{sec:met-limitations}

\paragraph{One evaluation set, twelve positives.} Every number in
Table~\ref{tab:met-selection} rests on $P = 12$. The exchange rate of
Equation~\eqref{eq:exchange} is exact given the class counts and does not
depend on sample size, but the empirical claim that the metric-preferred system
is the only one with signal rests on four MCC values whose sampling variability
is not characterised here. Chapter~\ref{ch:eval} reports a bootstrap interval
for ACCORD's AUROC, $[0.287, 0.633]$; no comparable interval is reported for any
MCC. [TODO: bootstrap all four MCCs on the $253$-pair evaluation set,
$B = 10{,}000$, and report intervals. If GPT-4o zero-shot's interval contains
zero, the claim that it is the only system with discriminative ability weakens
to a claim that it is the only system whose point estimate is positive.]

\paragraph{The confusion matrices are reconstructed, not read from an artifact.}
Table~\ref{tab:met-confusion} derives all sixteen cells from reported precision
and recall. The derivation is validated by reproducing four independently
reported MCC values to three decimals and by reproducing Chapter~\ref{ch:eval}'s
reported accuracy for ACCORD, which is strong evidence that it is correct, but
it remains a reconstruction. [TODO: emit the four confusion matrices directly
from the audit script and confirm cell by cell.]

\paragraph{Nine systems are missing.} The withdrawn submission compared thirteen
systems; four are recoverable at the level of traceability this thesis requires.
The chapter's conclusions do not depend on the missing nine, since the decisive
comparison is between the metric-preferred system and the highest-recall system
and both are present, but the comparison table is incomplete and is marked as
such.

\paragraph{No claim about field practice is made or supported.} This is a
deliberate retreat from the earlier submission. Whether papers in this area
select models by balanced accuracy or PABAK at low prevalence is an empirical
question about the literature, it was asserted without evidence, and it is
withdrawn. A survey establishing the base rates of metric reporting in clinical
NLP evaluations would settle it and was not conducted.

\paragraph{The severity-weighted analysis is not performed.}
Section~\ref{sec:met-prescription}, final paragraph. This is the most
substantive missing measurement in the chapter, because a triage system's value
is concentrated in the four contradictions rated $4$ or $5$ and none of the
statistics discussed here distinguishes them from the eight rated lower.

\section{Summary}
\label{sec:met-summary}

On an evaluation set with $P = 12$ positives and $Q = 241$ negatives, balanced
accuracy prefers a system that forgoes a true positive whenever it avoids more
than $Q/P = 20.08$ false positives; PABAK, which is $2\,\mathrm{accuracy} - 1$
and therefore not prevalence-robust in any useful sense, prefers it at a rate of
$1$; and MCC, which is not affine in the confusion cells, must be checked
numerically. Checked numerically, all three place a system recovering $6$ of
$12$ contradictions above two systems recovering $11$.

That much was the content of a withdrawn submission and it stands. What does not
stand is the framing. The systems recovering $11$ of $12$ do so by predicting
the positive class for over ninety percent of the evaluation set, at a
predicted-positive rate of exactly $1.000$ at every threshold from $0.05$ to
$0.15$, with AUROC $0.460$ and zero true contradictions in the top fifteen. Their
Matthews correlations are $-0.007$ and $-0.044$. The metric-preferred system, at
$+0.109$, is the only one in the comparison with discriminative ability. The
metrics did not select a worse system. They selected the only system there was,
out of a field in which the alternatives were noise.

The prescription that survives is narrower than the original claim and harder to
argue with: state $Q/P$ beside any balanced accuracy, do not present PABAK as
prevalence-robust, prespecify the primary measure and the cost ratio it encodes,
and report at least one statistic capable of returning the answer that no system
in the comparison has signal.
